\documentclass[11pt]{article}
\usepackage[margin=1in]{geometry}
\usepackage{amsmath,amssymb,hyperref}
\title{Projecting the prime exponents onto one field product}
\author{Project research note}
\date{18 November 2026}
\begin{document}
\maketitle

Write $b=\lceil\log_2 p\rceil$ and $B=\lfloor\sqrt p\rfloor$.
In the paired-slot setting of \texttt{exception\_product.tex}, the
exceptional correction over all terminal and chord factors is determined
by a single ordinary integer count once a primitive field element is
available. The construction has no atom or prime table and does not
query slot sizes. Its graph exceeds the budget of the optimized solver.
A second algorithm uses small auxiliary moduli and a certified bound
on rational height. Its cheap-query regime has a concrete limitation
on admissible hash windows.

\section{One count for the complete correction}
Let $F>0$ be the public integer cover. For a nonzero local weight $w$
use $h=\min(w,p-w)$ and $r=h/\gcd(h,F)$. On each component, the net
prime exponents $e_q$ from the Farey construction and its negative
sign count $C_-$ give
\[
 \prod r=(-1)^{C_-}\prod_{q\le B}q^{e_q}\pmod p .
\]
For a primitive $g$ the exponent of this product is the single
linear functional $(p-1)C_-/2+\sum_q e_q\log_g(q)$ modulo $p-1$.
Obtaining all coefficients $\log_g(q)$ in a public table would
already cost at least one entry per prime. Instead evaluate the
functional on an occurrence by certifying the unique $0\le e<p-1$
with $g^e=r$. This also folds in the sign of its Farey ratio.
The sign of the original centered weight belongs to the covered
factor and stays outside the correction.

Take a full row of $m$ points from the independently seeded symmetric
slots and require the lower canonical ordinate in slot one. Use a
single tag chosen from the terminal component and the chord
components $j=2,\ldots,m$.
\begin{itemize}
\item On the terminal tag, use $w_T(S_m)$, including the target and
infinity patches, for the running sum $S_m$.
\item On tag $j$ require a lower ordinate at slot $j$ and use the
nonzero chord scale $L(S_{j-1},P_j)$. Leave every later sign and slot
choice free. Each oriented prefix occurrence then has exactly
$\prod_{i>j}2n_i$ completions, its required continuation exponent.
\end{itemize}
For either tag add a canonical integer $0\le z<e$. The number of
extensions for the tagged row is $e$, including zero extensions for
$r=1$. All other auxiliaries are forced. If $Z$ is the total count,
then the complete correction is $g^{Z\bmod(p-1)}\bmod p$. A vacant slot
automatically gives $Z=0$. Counting all full rows with one tag also
avoids separate contractions and external repetitions for the prefixes.
To return the full field unit in one count, use $g^e=w$ in place of
$g^e=r$ on the same tags. This includes the covered factors and the
signs consumed by centering without a second query.

The graph uses the ordinary oriented curve addition with canonical
infinity coordinates and disjoint branches in \texttt{exception\_product.tex}.
It checks one prefix chain, including the unused suffix on a chord
tag. Canonical tag indicators and bitwise multiplexing across the
chain add $O(mb)$ edges. For example one can form local indicators
per tag, distribute them by bounded-degree equality trees, then
select each field bit by a binary tree with forced internal values.
The gcd check costs $O(b+\operatorname{bitlength}(F))$ as before.
The rank comparison is an ordinary $O(b)$-edge comparator. There
is no offset rounding with these deliberate witness multiplicities.
The exact count itself has at most
$O(b+\log m+\sum_i\log(2\tau_i))$ bits. A primitive $g$ can be
supplied with a certificate, or one can test uniform candidates by
counting $0<e<p-1$ with $g^e=1$. A zero count certifies primitivity.
The expected number of candidates is $(p-1)/\varphi(p-1)=O(b)$,
so this additional cost is polynomially many exponent graphs and
can be shared for a fixed field.

\paragraph{What the one query costs.}
In the existing block certificate in \texttt{unit\_recovery.tex}, use
$s$ exponent bits per block. Each block's $b$ output bits are split
among slice vertices. A slice reads the block digit and checks its
own bits of a fixed-base power by a polynomial-time local algorithm.
The number of Boolean edges is $\Theta(b\lceil b/s\rceil)$ for this
explicit-word layout, with additional $O(mb+\operatorname{bitlength}(F))$
for the rest of the graph. The largest degrees are
$O(s+b/s+\sqrt b\,\mathrm{polylog}\,b)$, including digit copies.
No table of $2^s$ entries is needed. At tuned $m$ and
$\operatorname{bitlength}(F)=O(b)$ the project's square-root-degree
contractor admits $s\asymp b^{2/3}$ and $S=O(b^{4/3})$ total edges.

The original external paper's Theorem 2 permits degree $O(S^{9/10})$,
polynomial description length and uniform polynomial-time local
predicates. Under that separate contraction statement we can take
$s\asymp b^{18/19}$. This gives $S=\Theta(b^{20/19})$ and degree
$O(b^{18/19})$. Equality subdivisions and the comparison with $z$
preserve the bounds. This is the optimal edge power within this
layout under that degree condition because $S=\Omega(b^2/s)$ and
$s=O(S^{9/10})$ imply $S=\Omega(b^{20/19})$. Theorem 2 only gives
$\log_2 T=O(b^{20/19}/\sqrt{\log b})$ for the resulting graph.
Its estimate is eventually larger in scale even than $b$.
Taking $s$ of order $b/m$ would keep $S=O(mb)$ but violate the
degree condition. These are layout and upper-bound comparisons,
not lower bounds for computing a field product.

\section{An alternative using small auxiliary fields}
In one component write its unsigned Farey atom product as the
reduced positive rational $N/D=\prod A_i/D_i$. Its field image
gets the separate Farey sign. Suppose $N,D\le 2^h$ with a certified
integer $h$. There is a way to recover this ratio using arithmetic
counts in small auxiliary characteristics. For an odd auxiliary
prime $\ell\ne p$ strip all powers of $\ell$ from each $A_i,D_i$.
Count the signed sum $v$ of the stripped valuations. If $v\ne0$,
skip this modulus. If $v=0$, the stripped-unit ratio gives $N/D$
modulo $\ell$, even when $\ell$ divided atoms on both sides.
Use a public generator and logarithm table in the small field to
sum the unit logs and exponentiate their difference.

For each atom its valuation is uniquely checked by
$A=\ell^j a$, $j\le\log_2 B$, $\ell\nmid a$. Select the public
power by the short index $j$ and enforce the quotient, remainder
and range checks. The low-width residue $a\bmod\ell$ determines
its small-field log by a local predicate. A bounded rank witness
counts each signed integer sum by subtracting two nonnegative
counts or using a known common shift. This also works with no
valuation factors present. For $\ell=\mathrm{poly}(b)$ the table
takes polynomial construction time and space and the added wires
are $O(b)$ per query. Complete oriented counts handle all path
multiplicities.

Accumulate good residues by CRT until their product exceeds
$2(2^h)^2$. A standard extended Euclidean reconstruction then
returns the unique positive reduced fraction of the promised
height. Two different candidates would have nonzero cross
difference of magnitude below the modulus, divisible by it.
For recovery, the approximating fraction
$(RD-N)/(MD)$ for a CRT representative $R$ and modulus $M$
reduces to $a/D$, where $a=(RD-N)/M$. It differs from $R/M$
by $N/(MD)<1/(2D^2)$, so it is a convergent, with the integral
case included. The implementation checks the recovered signs,
coprimality, congruence and both bounds.

The product of primes with $v\ne0$ divides $ND$. Their lost
logarithmic precision is at most $2h$. The Chebyshev lower bound
$\sum_{\ell\le x}\log\ell=\Omega(x)$ therefore guarantees enough
good moduli by $x=O(1+h)$, even after omitting 2 and $p$.
Constructing these fields, tables, CRT values and the reduced
fraction costs polynomial time in $b+h$. If $h=\mathrm{poly}(b)$
this gives polynomially many small-field queries with no full-width
field exponent certificate. In a general input $h$ can itself be
exponential in $b$.

\paragraph{How to certify the height.}
For the integer atom imbalance $c_d$ of the previous note, use the
exact three ordered-pair collision counts to compute
$J=\sum_{d\le B}c_d^2$. The paired graph has two independent
ordinary paths. Its arithmetic coefficient must be paid. Each of
the $M_0$ occurrences supplies one numerator and one denominator.
Thus $\sum c_d=0$ and its positive mass $P$ satisfies
\[
 P\le \min\{M_0,\lfloor J/2\rfloor,\lfloor\sqrt{BJ}/2\rfloor\}.
\]
Here $2P=\sum|c_d|\le J$ by integrality, and Cauchy--Schwarz
gives the other bound. Cancelling identical labels first bounds
both remaining integers by $B^P$. Reducing their gcd only decreases
them. Hence the displayed upper bound on $P$ times
$\lceil\log_2 B\rceil$ is a usable $h$. The population bound alone
also works without a pair count. A small number of agreeing
auxiliary residues does not certify height. For instance $1$ and
$2^{60}$ agree modulo $3,5,7,11$ but are $1,123$ modulo 211.

There are admissible one-slot seeds on which the true rational
height is already $\Omega(\sqrt p\log p)$ bits. Take the populated
interval of width $W\asymp\sqrt p$ with $W\le\sqrt p/2$ from
\texttt{exception\_product.tex} and set $F=1$. Its positive
unpatched distances give Farey ratios $d/1$. There are
$\Omega(W)$ distinct distances, each realized at least once, so
the integer product is at least $\lfloor cW\rfloor!$ for a fixed
$c>0$ on the stated curve class. The denominator is one. This
explains the limit of reconstructing the full rational before field
reduction. It is a family of allowable windows and does not prove a
lower bound for independent tuned multi-slot seeds or other algorithms.

\section{Comparison with the external paper}
Akshima, Bacho, Huang, Kiltz, Nishaburi, Nooripoor, Tittelbach and
Wiedijk, \emph{A Subexponential Time Algorithm for Elliptic Curve
Discrete Logarithm}, version 0.2, 29 October 2026, Theorems 1, 2
and 4, is the external baseline.\footnote{\url{https://github.com/bochum-crypto/subexponential-ecdl/blob/main/subexponential-ecdl.pdf}}
For prime-order curves over $b$-bit prime fields their expected
time is $2^{O(b/(\log_2 b)^{1/4})}$ conditional on the cited
Boolean tensor contraction theorem. They also discuss fixed cofactors.
Their main ECDL theorem does not state a numerical leading time
coefficient on this scale. A percentage improvement against that
theorem cannot be inferred.

For an explicit numerical comparison inside the project, the
supplied compressed-relation package reports $27.0138774237$.
The short inverse-hash circuit reduces the leading word budget
from $38mb$ to $36mb$ and the time coefficient to $26.2933775156$,
a $1-\sqrt{18/19}=2.6671\%$ reduction on the common input class.
The supplied simplex manuscript at project intake reported
$32.1569516597$ for the common supplied-order polynomial-cofactor
class. The current time coefficient is $18.2342\%$ below that earlier
project baseline. The conditional workspace coefficient is $13.1466887578$.
These constants use the local finite-packet Assumption A.1 and
the inherited arithmetic, contraction and relation-reduction
results. Neither the old nor the new scalar projection changes
the optimizing balance between $b/m$ and $mb/\sqrt{\log b}$.
The intended $2^{\widetilde O(\sqrt b)}$ bound and an end-to-end
speed demonstration on large inputs remain open.

\section{Verification}
Run \texttt{python -B research/single\_target/verify\_prime\_projection.py}.
The two reductions agree with independently formed component products
on 24 small seed/cover cases. For the 78-exception fixture the
single count is 34,437 with $g=2$, yielding the complete correction
$2^{34437}\bmod 211=132$, including the chord factor. The full-weight
variant independently gives 93 and agrees with the literal normalized
polynomial product on all twelve small seed cases.
The verifier also checks 2,653 bounded rational reconstructions.
On the twelve four-slot fixtures over $\mathbb F_{65521}$ all tagged
counts give the expected complete correction. The exact terminal
heights $\lceil\log_2\max(N,D)\rceil$ range from 2,205 to 14,349
bits, and the computable collision bounds range from 13,256 to
94,432 bits. These heights refer to the positive rational before
field reduction. The finite oracles enumerate tuples. The larger
scalar test uses a field-sized reference log table, which is not
assumed by the graph construction or included as evidence of a
contractor speedup.
\end{document}
